% Part: counterfactuals
% Chapter: introduction
% Section: paradoxes-material

\documentclass[../../../include/open-logic-section]{subfiles}

\begin{document}

\olfileid{cnt}{int}{par}

\olsection{پارادوکس‌های شرطی مادی}

سی.~آی. لوئیس از نخستین کسانی بود که کاربرد~$!A \lif !B$ را برای
نمادگذاری عبارت‌های انگلیسی «اگر \dots، آنگاه \dots» نقد کرد. لوئیس
به کاربرد شرطی مادی در \emph{اصول ریاضیات} وایتهد و راسل اعتراض داشت؛
آنان~$\lif$ را «مستلزم است» می‌خواندند. لوئیس به‌درستی اعتراض کرد که
اگر~$\lif$ به معنای «مستلزم است» باشد، هر گزارهٔ کاذب~$p$ مستلزم آن
است که~$p$ مستلزم~$q$ باشد، زیرا اگر~$p$ کاذب باشد
$p \lif (p \lif q)$ صادق است؛ و هر گزارهٔ صادق~$q$ مستلزم آن است که
$p$ مستلزم~$q$ باشد، زیرا اگر~$q$ صادق باشد~$q \lif (p \lif q)$ صادق
است.

البته منطق‌دانان می‌دانند \emph{استلزام}، یعنی استلزام منطقی، یک رابط
نیست بلکه رابطه‌ای میان !!{formula}ها یا گزاره‌هاست. پس برای پرهیز از
خلط، کافی است~$\lif$ را «مستلزم است» نخوانیم.\footnote{خواندن~«$\lif$»
  به‌صورت «مستلزم است» هنوز میان ریاضی‌دانان و دانشمندان رایانه بسیار
  رایج است، هرچند فیلسوفان برای پرهیز از خلط‌هایی که لوئیس برجسته کرد،
  می‌کوشند آن را «تنها اگر» بخوانند.} مادام که چنین نکنیم، نگرانی خاص
لوئیس اصلاً پیش نمی‌آید: حتی اگر~$p$ را نمایندهٔ جمله‌ای کاذب در
انگلیسی بدانیم، $p$، $q$ را «نتیجه نمی‌دهد». برای تعیین اینکه آیا
$p \Entails q$، باید \emph{همهٔ} !!{valuation}ها را در نظر بگیریم، و
حتی وقتی~$p$ را برای نمادگذاری جمله‌ای به کار می‌بریم که اتفاقاً کاذب
است، $p \Entails/ q$.

بااین‌حال، هنوز در عبارت‌های «اگر \dots، آنگاه \dots» مانند مثال لوئیس
چیزی نامتعارف هست:
\begin{quote}
اگر ماه از پنیر سبز ساخته شده باشد، آنگاه~$2+2=4$.
\end{quote}
و نیز در استنتاج‌های زیر:
\begin{quote}
  ماه از پنیر سبز ساخته نشده است. بنابراین، اگر ماه از پنیر سبز ساخته
  شده باشد، آنگاه~$2+2=4$.

  $2+2 = 4$. بنابراین، اگر ماه از پنیر سبز ساخته شده باشد، آنگاه
  $2+2=4$.
\end{quote}
اما اگر «اگر \dots، آنگاه \dots» صرفاً همان~$\lif$ می‌بود، جمله
بی‌هیچ اشکالی صادق و استنتاج‌ها بی‌هیچ اشکالی معتبر می‌بودند.

نمونه‌ای دیگر به همان‌گویی~$(!A \lif !B) \lor (!B \lif !A)$ مربوط
است. این امر القا می‌کند که اگر دو جملهٔ اخباری~$S$ و~$T$ را تصادفی از
روزنامه بردارید، جملهٔ «اگر~$S$ آنگاه~$T$، یا اگر~$T$ آنگاه~$S$» باید
صادق باشد.

\end{document}
